\section{Round-5 Objective}
\label{sec:round5-objective}
\textbf{Path pursued: (C) obstruction/correction, by settling the upper-construction side rigorously.}

Round~4 isolated the main unresolved issue as the behaviour near the exponent $1/2$: it produced structural obstructions on any putative counterexample to the lower-bound direction, but it did not rigorously prove the matching upper construction $|[1,N]\setminus B|\ll_\varepsilon N^{1/2+\varepsilon}$.

In this round I prove exactly that upper-construction statement. Thus the remaining open part of Erd\H{o}s Problem~\#14 is not the existence of near-optimal constructions from above, but whether one can beat the square-root threshold from below, i.e. whether $o(N^{1/2})$ is possible or whether a lower bound $\gg_\varepsilon N^{1/2-\varepsilon}$ always holds.

\section{Round-4 Foundation Used}
\label{sec:round4-foundation-used}
I use the following Round~4 conclusions.
\begin{itemize}
\item The Round~4 global density obstruction showed that if $E_N:=|[1,N]\setminus B|$ were as small as $o(N^{1/2})$, then the counting function of $A$ would have to lie in a narrow window, roughly between $N^{1/2}$ and $N^{3/4}$.
\item More importantly for the present round, Round~4 left open the exact construction side: produce a single infinite set $A\subseteq \mathbb N$ with
\[
E_N\ll_\varepsilon N^{1/2+\varepsilon}.
\]
\end{itemize}
No Round~4 lemma is needed inside the new proof; Round~4 is used to identify the exact remaining gap and to interpret the significance of the new theorem.

\section{New Insight / Tool}
\label{sec:round5-new-tool}
The new tool is a \emph{balanced digit-splitting construction in base $s^2$}.

For an integer $s\ge 2$, let
\[
q:=s^2.
\]
I construct two sets $F_s,G_s\subseteq \mathbb N_0$ such that:
\begin{enumerate}
\item every nonnegative integer has a \emph{unique} representation as $f+g$ with $f\in F_s$ and $g\in G_s$;
\item the self-sum sets $F_s+F_s$ and $G_s+G_s$ have counting function
\[
\ll_s N^{\beta_s},\qquad \beta_s:=\frac{\log(2s-1)}{\log(s^2)}=\frac{\log(2s-1)}{2\log s};
\]
\item since $\beta_s\downarrow 1/2$ as $s\to\infty$, this yields, for every $\varepsilon>0$, an infinite set $A\subseteq\mathbb N$ with
\[
|[1,N]\setminus B|\ll_\varepsilon N^{1/2+\varepsilon}.
\]
\end{enumerate}
The case $s=2$ recovers the classical exponent $\log 3/\log 4$; the new point is that allowing arbitrary $s$ pushes the exponent down to $1/2+\varepsilon$.

\section{Attack Plan}
\label{sec:round5-attack-plan}
The exact Round~4 gap is the absence of a rigorous infinite construction matching the heuristic square-root upper scale.

The plan is:
\begin{enumerate}
\item define $F_s$ and $G_s$ by digit restrictions in base $q=s^2$;
\item prove that every $m\in\mathbb N_0$ has a unique decomposition $m=f+g$ with $f\in F_s$, $g\in G_s$;
\item prove that $F_s+F_s$ and $G_s+G_s$ are sparse, with exponent $\beta_s$;
\item define
\[
A_s:=(F_s+1)\cup(G_s+1)\subseteq\mathbb N
\]
and show that every $n\ge 2$ has one canonical representation coming from $F_s+G_s$, while every failure of uniqueness is forced into the sparse exceptional sets $(F_s+F_s)+2$ or $(G_s+G_s)+2$.
\end{enumerate}
This gives a full, gap-free proof of the upper-construction statement.

\section{Work}
\label{sec:round5-work}
For a set $X\subseteq\mathbb N_0$ and $N\ge 0$, write
\[
X(N):=|X\cap [0,N]|.
\]
For $A\subseteq\mathbb N$, as before,
\[
r_A(n):=|\{\{a,b\}:a,b\in A,\ a\le b,\ a+b=n\}|,
\qquad
B(A):=\{n\in\mathbb N:r_A(n)=1\}.
\]

\subsection*{Lemma 5.1 (digit splitting)}
Fix an integer $s\ge 2$ and put $q=s^2$.
Define
\[
F_s:=\left\{\sum_{j\ge 0} u_j q^j:
0\le u_j\le s-1\ \text{for all $j$, and $u_j=0$ for all but finitely many $j$}\right\},
\]
and
\[
G_s:=sF_s=\left\{\sum_{j\ge 0} sv_j q^j:
0\le v_j\le s-1\ \text{for all $j$, and $v_j=0$ for all but finitely many $j$}\right\}.
\]
Then every integer $m\ge 0$ has a unique representation
\[
m=f+g,\qquad f\in F_s,\ g\in G_s.
\]
Moreover,
\[
F_s\cap G_s=\{0\}.
\]

\paragraph{Proof.}
Write the base-$q$ expansion of $m$ as
\[
m=\sum_{j\ge 0} d_j q^j,\qquad 0\le d_j\le q-1=s^2-1.
\]
For each $j$, divide $d_j$ by $s$:
\[
d_j=u_j+sv_j,\qquad 0\le u_j\le s-1,\quad 0\le v_j\le s-1.
\]
This choice is unique. Therefore
\[
f:=\sum_{j\ge 0} u_j q^j\in F_s,
\qquad
g:=\sum_{j\ge 0} sv_j q^j\in G_s,
\]
and $m=f+g$.

Uniqueness is digitwise: if $m=f+g=f'+g'$ with $f,f'\in F_s$ and $g,g'\in G_s$, then comparing base-$q$ digits yields the same quotient-and-remainder decomposition $d_j=u_j+sv_j$, so $f=f'$ and $g=g'$.

Finally, if $n\in F_s\cap G_s$, then each base-$q$ digit of $n$ lies both in $\{0,1,\dots,s-1\}$ and in $\{0,s,2s,\dots,(s-1)s\}$. Their intersection is $\{0\}$, so every digit is $0$ and hence $n=0$.
\hfill$\square$

\subsection*{Lemma 5.2 (sparse self-sums)}
With $F_s,G_s$ as above, define
\[
\beta_s:=\frac{\log(2s-1)}{\log(s^2)}=\frac{\log(2s-1)}{2\log s}.
\]
Then
\[
(F_s+F_s)(N)\ll_s N^{\beta_s}
\qquad\text{and}\qquad
(G_s+G_s)(N)\ll_s N^{\beta_s}.
\]

\paragraph{Proof.}
Take $f,f'\in F_s$, say
\[
f=\sum_{j\ge 0} u_j q^j,
\qquad
f'=\sum_{j\ge 0} u'_j q^j,
\qquad 0\le u_j,u'_j\le s-1.
\]
Then
\[
f+f'=\sum_{j\ge 0}(u_j+u'_j)q^j,
\qquad 0\le u_j+u'_j\le 2s-2<q.
\]
So there are \emph{no carries}. Hence $F_s+F_s$ consists exactly of those integers whose base-$q$ digits all lie in $\{0,1,\dots,2s-2\}$.

If $q^m\le N<q^{m+1}$, then every element of $(F_s+F_s)\cap[0,N]$ is determined by at most $m+1$ base-$q$ digits, each having $2s-1$ possibilities. Therefore
\[
(F_s+F_s)(N)\le (2s-1)^{m+1}
\ll_s N^{\log_q(2s-1)}
= N^{\beta_s}.
\]

Since $G_s=sF_s$, we have
\[
G_s+G_s=s(F_s+F_s),
\]
so
\[
(G_s+G_s)(N)\le (F_s+F_s)(N/s)\ll_s N^{\beta_s}.
\]
This proves the lemma.
\hfill$\square$

\subsection*{Proposition 5.3 (explicit infinite construction)}
For each integer $s\ge 2$, set
\[
A_s:=(F_s+1)\cup(G_s+1)\subseteq\mathbb N.
\]
Then $A_s$ is infinite and
\begin{equation}
\label{eq:As-upper-bound}
|[1,N]\setminus B(A_s)|\ll_s N^{\beta_s},
\qquad \beta_s=\frac{\log(2s-1)}{2\log s}.
\end{equation}

\paragraph{Proof.}
First, $A_s$ is infinite because $F_s$ is infinite.

Let $n\ge 2$. By Lemma~5.1, $n-2$ has a unique representation
\[
n-2=f+g,\qquad f\in F_s,\ g\in G_s.
\]
Therefore
\[
n=(f+1)+(g+1),
\]
so every $n\ge 2$ is represented as a sum of two elements of $A_s$.

Now suppose that $n\notin B(A_s)$. Then either $n=1$, or $n\ge 2$ and there is some representation
\[
n=(x+1)+(y+1),\qquad x,y\in F_s\cup G_s,
\]
which is different from the canonical unordered pair $\{f+1,g+1\}$ coming from $n-2=f+g$.

If one of $x,y$ lies in $F_s$ and the other lies in $G_s$, then $x+y=n-2$ is a mixed decomposition of $n-2$. By the uniqueness in Lemma~5.1, this forces
\[
\{x,y\}=\{f,g\},
\]
so it gives the same unordered representation of $n$ and does \emph{not} create a second one.

Hence any genuine failure of uniqueness must come from a same-type representation: either
\[
x,y\in F_s
\qquad\text{or}\qquad
x,y\in G_s.
\]
Consequently,
\[
[1,N]\setminus B(A_s)
\subseteq
\{1\}\cup \bigl((F_s+F_s)+2\bigr)\cup\bigl((G_s+G_s)+2\bigr).
\]
Therefore, by Lemma~5.2,
\[
|[1,N]\setminus B(A_s)|
\le 1+(F_s+F_s)(N-2)+(G_s+G_s)(N-2)
\ll_s N^{\beta_s}.
\]
This proves \eqref{eq:As-upper-bound}.
\hfill$\square$

\subsection*{Corollary 5.4 (the $N^{1/2+\varepsilon}$ construction)}
For every $\varepsilon>0$ there exists an infinite set $A\subseteq\mathbb N$ such that
\[
|[1,N]\setminus B(A)|\ll_{\varepsilon} N^{1/2+\varepsilon}.
\]

\paragraph{Proof.}
For $s\ge 2$,
\[
\beta_s
=\frac{\log(2s-1)}{2\log s}
=\frac12+\frac{\log(2-1/s)}{2\log s}
<\frac12+\frac{\log 2}{2\log s}.
\]
Choose $s$ so large that $\beta_s<\tfrac12+\varepsilon$; for instance any
\[
s>2^{1/(2\varepsilon)}
\]
suffices. Then Proposition~5.3 gives an infinite set $A_s\subseteq\mathbb N$ with
\[
|[1,N]\setminus B(A_s)|\ll_s N^{\beta_s}\ll_{\varepsilon} N^{1/2+\varepsilon}.
\]
\hfill$\square$

\subsection*{Consequence for the original problem}
This round proves the full upper-construction side that had remained missing in the previous rounds:
\[
\forall\varepsilon>0\ \exists\, A\subseteq\mathbb N\ \text{infinite such that}\ |[1,N]\setminus B(A)|\ll_{\varepsilon}N^{1/2+\varepsilon}.
\]
Therefore the unresolved part of Erd\H{o}s Problem~\#14 is now concentrated entirely on the \emph{lower} side: whether every infinite $A\subseteq\mathbb N$ must satisfy
\[
|[1,N]\setminus B(A)|\gg_{\varepsilon}N^{1/2-\varepsilon}
\]
for all large $N$, or whether one can in fact achieve $o(N^{1/2})$.

\section{Adversarial Verification}
\label{sec:round5-adversarial}
I check the new construction against the main possible failure modes.
\begin{itemize}
\item \textbf{Carry issue in the self-sums.} For $F_s+F_s$, the largest digit sum is $2s-2$, and since $2s-2<s^2=q$, there are no carries. This is the key point making the exponent $\beta_s$ exact. For $G_s+G_s$, I avoid a separate carry analysis by using the identity $G_s=sF_s$.

\item \textbf{Overlap of the two parts.} We proved $F_s\cap G_s=\{0\}$. After shifting by $1$, the two parts intersect only at $1$. This does not create extra unordered representations: if a mixed representation of $n-2$ exists, Lemma~5.1 makes it unique, so it contributes exactly one unordered pair.

\item \textbf{Possibility of missing sums.} There are no missing sums for $n\ge 2$, because every $n-2\in\mathbb N_0$ has a mixed representation $f+g$. The integer $1$ is the only inevitable missing sum in the shifted positive setting.

\item \textbf{Boundary case $n=2$.} Here $n-2=0=0+0$, so the canonical representation is $2=1+1$. This is still exactly one unordered representation.

\item \textbf{Dependence on $s$.} All constants in Proposition~5.3 depend only on the fixed parameter $s$. After choosing $s=s(\varepsilon)$, this becomes the claimed $\ll_\varepsilon$ bound.
\end{itemize}
No hidden assumption beyond the adopted unordered-representation convention is used.

\section{Final}
\label{sec:round5-final}
\textbf{UNRESOLVED (BUT STRICTLY ADVANCED).}

This round completely settles the \emph{upper-construction side} left open in the previous rounds:
\[
\forall\varepsilon>0\ \exists\, A\subseteq\mathbb N\ \text{infinite such that}\ |[1,N]\setminus B(A)|\ll_{\varepsilon}N^{1/2+\varepsilon}.
\]
The proof is explicit and gap-free.

So the surviving open core of Erd\H{o}s Problem~\#14 is now precisely this:
\begin{quote}
Can one prove a matching lower bound of order $N^{1/2-o(1)}$ for every infinite $A$, or can one construct a single infinite $A$ with $|[1,N]\setminus B(A)|=o(N^{1/2})$?
\end{quote}
In other words, the exponent $1/2$ is now rigorously established as the correct \emph{upper} threshold, and only the lower-threshold question remains open.

\section{Completion Estimate}
\label{sec:round5-completion}
\textbf{COMPLETION: 84\%}

\section{References}
\label{sec:round5-references}
\begin{itemize}
\item Internal previous-round files: \texttt{14\_v2.tex}, \texttt{14\_v2\_solutions.tex}, \texttt{14\_solutions.tex}.
\item T. F. Bloom, \emph{Erd\H{o}s Problem \#14}, \texttt{erdosproblems.com/14}, accessed 2026-03-07.
\item A. S\'ark\"ozy and V. T. S\H{o}s, \emph{On Additive Representation Functions}, in \emph{The Mathematics of Paul Erd\H{o}s I}, Springer, 1997.
\end{itemize}
